% 本文件由 tools/gen_terms.py 自动生成，请勿手工编辑。
% 数据源：deliverable/术语对照表.tsv
\chapter*{术语对照表}
\addcontentsline{toc}{chapter}{术语对照表}
\markboth{术语对照表}{}
\noindent 下表列出本书主要术语的英文原词与中译名。「处理」一列：\textbf{译} 表示正文用中文译名；\textbf{保留} 表示正文保留拉丁原文（多为专名或通行缩写）；\textbf{首见} 表示首次出现时给出「中文（English）」后文用中文。人名一律保留原文，不列于此。\par
\medskip

\newcommand{\termrowA}[2]{\par\noindent\interlinepenalty=10000\hangindent=1.1em\hangafter=1 \textbf{#1}\quad #2\par}
\newcommand{\termrowB}[3]{\par\noindent\interlinepenalty=10000\hangindent=1.1em\hangafter=1 \textbf{#1}\quad #2\par\nobreak\hspace*{1.1em}{\small\color{black!65}\interlinepenalty=10000 #3\par}}
\newcommand{\termgroup}[1]{\par\medskip\noindent{\bfseries\large #1}\par\medskip}

\begin{multicols}{2}
\raggedright
\setlength{\parskip}{0.35em}

\termgroup{一、点集拓扑与流形}
\termrowA{topological space}{拓扑空间}
\termrowA{topology}{拓扑}
\termrowA{open set}{开集}
\termrowA{closed set}{闭集}
\termrowA{neighborhood}{邻域}
\termrowA{interior}{内部}
\termrowA{closure}{闭包}
\termrowB{boundary}{边界}{链的 boundary 亦作「边缘」，见「同调」组}
\termrowB{compact}{紧}{「紧致」弃用}
\termrowA{locally compact}{局部紧}
\termrowB{Hausdorff}{Hausdorff}{〔保留〕 不加「豪斯多夫」}
\termrowA{second countable}{第二可数}
\termrowA{paracompact}{仿紧}
\termrowA{metric space}{度量空间}
\termrowA{continuous map}{连续映射}
\termrowA{bijective}{双射}
\termrowA{injective / surjective}{单射 / 满射}
\termrowA{subspace}{子空间}
\termrowA{quotient space}{商空间}
\termrowA{quotient map}{商映射}
\termrowA{equivalence relation}{等价关系}
\termrowA{homeomorphism}{同胚}
\termrowA{homeomorphic}{同胚的}
\termrowA{topological manifold}{拓扑流形}
\termrowA{smooth manifold}{光滑流形}
\termrowA{dimension}{维数}
\termrowB{coordinate patch}{坐标卡}{底本常用 patches，复数「坐标卡」}
\termrowA{chart}{卡}
\termrowA{atlas}{图册}
\termrowA{transition function}{转移函数}
\termrowA{manifold with boundary}{带边流形}
\termrowA{orientable}{可定向}
\termrowA{non-orientable}{不可定向}
\termrowA{orientation}{定向}
\termrowA{connected sum}{连通和}
\termrowA{genus}{亏格}
\termrowA{sphere}{球面}
\termrowA{torus}{环面}
\termrowB{real / complex projective space}{实 / 复射影空间}{记号 $\RP^n$、$\CP^n$}
\termrowB{Möbius band}{Möbius 带}{〔保留〕}
\termrowB{Klein bottle}{Klein 瓶}{〔保留〕}
\termrowB{Euler characteristic}{Euler 示性数}{〔保留〕}
\termrowA{surface}{曲面}
\termrowA{embedding}{嵌入}
\termrowA{immersion}{浸入}
\termrowA{submanifold}{子流形}
\termrowA{triangulation}{三角剖分}
\termrowB{simplex}{单形}{复数 simplexes / simplices 均作「单形」}
\termrowA{polyhedron}{多面体}
\termrowA{cell}{胞腔}
\termrowA{CW complex}{CW 复形}
\termrowA{skeleton}{骨架}
\termrowB{covering space}{覆叠空间}{「覆盖空间」弃用，全书统一「覆叠」}
\termrowA{universal cover}{万有覆叠}
\termrowA{deck transformation}{覆叠变换}
\termrowA{lift}{提升}
\termrowA{Lie group}{李群}
\termrowA{Lie algebra}{李代数}
\termrowA{homogeneous space}{齐性空间}
\termrowA{orbit}{轨道}
\termrowA{stabilizer}{稳定化子}
\termrowA{free action}{自由作用}
\termrowA{transitive}{传递的}
\termrowA{coset}{陪集}
\termrowB{Grassmannian}{Grassmann 流形}{〔保留〕}
\termrowA{projective space}{射影空间}
\termrowA{connected}{连通}
\termrowA{simply connected}{单连通}
\termrowA{contractible}{可缩的}
\termrowA{deformation retract}{形变收缩}
\termrowA{homotopy equivalence}{同伦等价}

\termgroup{二、纤维丛、联络与规范场}
\termrowA{fiber bundle}{纤维丛}
\termrowA{fiber}{纤维}
\termrowA{base space}{底空间}
\termrowA{total space}{全空间}
\termrowA{projection}{投影}
\termrowA{local trivialization}{局部平凡化}
\termrowA{structure group}{结构群}
\termrowB{principal G-bundle}{主 $G$-丛}{首见写「主 $G$-丛（principal $G$-bundle）」}
\termrowA{associated bundle}{相伴丛}
\termrowA{vector bundle}{向量丛}
\termrowA{line bundle}{线丛}
\termrowA{tangent bundle}{切丛}
\termrowA{cotangent bundle}{余切丛}
\termrowA{normal bundle}{法丛}
\termrowA{frame bundle}{标架丛}
\termrowA{universal bundle}{万有丛}
\termrowA{tautological bundle}{重言丛}
\termrowA{pull-back bundle}{拉回丛}
\termrowA{section}{截面}
\termrowA{fiber bundle map}{丛映射}
\termrowA{bundle isomorphism}{丛同构}
\termrowA{classifying space}{分类空间}
\termrowA{classifying map}{分类映射}
\termrowA{tangent space}{切空间}
\termrowA{cotangent space}{余切空间}
\termrowA{vector field}{向量场}
\termrowA{differential form}{微分形式}
\termrowA{exterior derivative}{外微分}
\termrowA{wedge product}{楔积}
\termrowA{pull-back}{拉回}
\termrowA{push-forward}{前推}
\termrowA{interior product}{内乘}
\termrowA{closed form}{闭形式}
\termrowA{exact form}{恰当形式}
\termrowA{volume form}{体积形式}
\termrowA{integration}{积分}
\termrowB{Stokes' theorem}{Stokes 定理}{〔保留〕}
\termrowB{de Rham cohomology}{de Rham 上同调}{〔保留〕}
\termrowB{de Rham complex}{de Rham 复形}{〔保留〕}
\termrowA{connection}{联络}
\termrowA{curvature}{曲率}
\termrowA{gauge field}{规范场}
\termrowA{gauge potential}{规范势}
\termrowA{field strength}{场强}
\termrowA{gauge group}{规范群}
\termrowA{gauge theory}{规范理论}
\termrowA{gauge transformation}{规范变换}
\termrowB{Maxwell field}{Maxwell 场}{〔保留〕}
\termrowB{Yang--Mills}{Yang--Mills}{〔保留〕}
\termrowB{electromagnetism}{电磁相互作用}{单说 electromagnetism 时作「电磁现象」}
\termrowA{strong force}{强相互作用}
\termrowA{electromagnetic}{电磁的}
\termrowB{monopole}{磁单极子}{Dirac monopole「Dirac 磁单极子」}
\termrowA{instanton}{瞬子}
\termrowA{soliton}{孤子}
\termrowA{topological soliton}{拓扑孤子}
\termrowA{Standard Model}{标准模型}
\termrowB{spin structure}{spin 结构}{〔保留〕 不译「自旋结构」}
\termrowB{spin manifold}{spin 流形}{〔保留〕}
\termrowB{spin connection}{spin 联络}{〔保留〕}
\termrowA{theta term}{$\theta$ 项}
\termrowB{Chern--Simons term}{Chern--Simons 项}{〔保留〕}
\termrowB{Wess--Zumino--Witten (WZW) term}{WZW 项}{〔保留〕 首见写「Wess--Zumino--Witten（WZW）项」}

\termgroup{三、同伦}
\termrowA{homotopy}{同伦}
\termrowA{homotopic}{同伦的}
\termrowA{homotopy group}{同伦群}
\termrowA{fundamental group}{基本群}
\termrowA{higher homotopy group}{高阶同伦群}
\termrowA{base point}{基点}
\termrowA{pointed space}{带基点的空间}
\termrowA{path}{道路}
\termrowA{loop}{环路}
\termrowA{loop space}{环路空间}
\termrowA{path space}{道路空间}
\termrowA{concatenation}{连接}
\termrowA{abelian}{阿贝尔的}
\termrowA{commutative}{交换的}
\termrowA{fibration}{纤维化}
\termrowA{exact sequence}{正合列}
\termrowA{short exact sequence}{短正合列}
\termrowA{long exact sequence}{长正合列}
\termrowA{connecting homomorphism}{连接同态}
\termrowA{suspension}{悬挂}
\termrowA{suspension isomorphism}{悬挂同构}
\termrowB{Eilenberg--MacLane space}{Eilenberg--MacLane 空间}{〔保留〕}
\termrowB{Postnikov tower}{Postnikov 塔}{〔保留〕}
\termrowB{Whitehead theorem}{Whitehead 定理}{〔保留〕}
\termrowA{weak homotopy equivalence}{弱同伦等价}
\termrowA{homotopy fiber}{同伦纤维}
\termrowA{classifying space BG}{分类空间 $BG$}
\termrowA{universal cover}{万有覆叠}
\termrowB{Seifert--van Kampen?}{Seifert--van Kampen 定理}{〔保留〕}
\termrowB{$\pi_n(S^m)$}{$\pi_n(S^m)$}{〔保留〕}
\termrowA{degree (of a map)}{映射的度}
\termrowA{winding number}{缠绕数}
\termrowA{topological defect}{拓扑缺陷}
\termrowA{defect}{缺陷}

\termgroup{四、同调与上同调}
\termrowA{homology}{同调}
\termrowA{cohomology}{上同调}
\termrowA{(singular) chain}{（奇异）链}
\termrowA{chain complex}{链复形}
\termrowA{boundary operator}{边缘算子}
\termrowA{cycle}{闭链}
\termrowB{boundary (of a chain)}{边缘}{与「边界 boundary of a manifold」区分：链的边缘}
\termrowA{homologous}{同调的}
\termrowA{homology class}{同调类}
\termrowA{cohomology class}{上同调类}
\termrowA{cochain}{上链}
\termrowA{cocycle}{上闭链}
\termrowA{coboundary}{上边缘}
\termrowA{singular homology}{奇异同调}
\termrowA{singular cohomology}{奇异上同调}
\termrowA{simplicial homology}{单纯同调}
\termrowA{simplicial cohomology}{单纯上同调}
\termrowA{cellular homology}{胞腔同调}
\termrowA{cellular cohomology}{胞腔上同调}
\termrowB{Čech cohomology}{Čech 上同调}{〔保留〕}
\termrowA{coefficient group}{系数群}
\termrowB{coefficients in $A$}{以 $A$ 为系数}{写「$H^n(M;A)$，以 $A$ 为系数」}
\termrowB{torsion}{挠}{torsion subgroup「挠子群」}
\termrowA{free part}{自由部分}
\termrowB{Betti number}{Betti 数}{〔保留〕}
\termrowB{Euler characteristic}{Euler 示性数}{〔保留〕}
\termrowA{universal coefficient theorem}{万有系数定理}
\termrowB{Künneth formula}{Künneth 公式}{〔保留〕}
\termrowB{Hurewicz theorem}{Hurewicz 定理}{〔保留〕}
\termrowA{cup product}{上积}
\termrowA{cap product}{下积}
\termrowA{cross product}{叉积}
\termrowA{relative homology}{相对同调}
\termrowA{excision}{切除}
\termrowB{Mayer--Vietoris sequence}{Mayer--Vietoris 序列}{〔保留〕}
\termrowA{homotopy invariance}{同伦不变性}
\termrowA{fundamental class}{基本类}
\termrowA{orientation class}{定向类}
\termrowB{Thom class}{Thom 类}{〔保留〕}
\termrowB{Poincaré duality}{Poincaré 对偶}{〔保留〕}
\termrowB{toric code}{环面码}{首见写「环面码（toric code）」}
\termrowA{stabilizer code}{稳定子码}
\termrowA{lattice gauge theory}{格点规范理论}

\termgroup{五、示性类、指标与自旋}
\termrowA{characteristic class}{示性类}
\termrowB{Chern class}{Chern 类}{〔保留〕}
\termrowB{Stiefel--Whitney class}{Stiefel--Whitney 类}{〔保留〕 记号 $w_i$}
\termrowB{Pontryagin class}{Pontryagin 类}{〔保留〕 记号 $p_i$}
\termrowB{Euler class}{Euler 类}{〔保留〕}
\termrowB{Dixmier--Douady class}{Dixmier--Douady 类}{〔保留〕}
\termrowB{first Chern class}{第 1 Chern 类}{〔保留〕 底本用「1st Chern class」，中文写「第 1 Chern 类」}
\termrowB{Chern--Weil theory}{Chern--Weil 理论}{〔保留〕}
\termrowB{integral cohomology}{整系数上同调}{integral class「整的类」}
\termrowA{charactersitic class of a bundle}{丛的示性类}
\termrowB{spin structure}{spin 结构}{〔保留〕}
\termrowB{spin$^c$ structure}{spin$^c$ 结构}{〔保留〕}
\termrowB{Pin group}{Pin 群}{〔保留〕}
\termrowB{pinor}{pinor}{〔保留〕 不译}
\termrowA{spinor}{旋量}
\termrowB{Clifford algebra}{Clifford 代数}{〔保留〕}
\termrowA{gamma matrix}{$\gamma$ 矩阵}
\termrowA{chirality}{手征性}
\termrowA{chiral}{手征的}
\termrowB{Dirac operator}{Dirac 算子}{〔保留〕}
\termrowA{index theorem}{指标定理}
\termrowB{Atiyah--Singer index theorem}{Atiyah--Singer 指标定理}{〔保留〕}
\termrowB{eta invariant}{eta 不变量}{〔保留〕}
\termrowA{spectral flow}{谱流}
\termrowB{Atiyah--Patodi--Singer}{Atiyah--Patodi--Singer}{〔保留〕 缩写 APS}
\termrowB{Pfaffian}{Pfaffian}{〔保留〕}
\termrowA{anomaly}{反常}
\termrowA{chiral anomaly}{手征反常}
\termrowA{parity anomaly}{宇称反常}
\termrowA{anomaly inflow}{反常流入}
\termrowA{anomaly polynomial}{反常多项式}
\termrowB{level (Chern--Simons level)}{级}{Chern--Simons 的 level $k$ 作「级 $k$」}
\termrowA{Hall conductance}{霍尔电导}
\termrowA{quantum Hall effect}{量子霍尔效应}
\termrowA{charge quantization}{电荷量子化}
\termrowB{Wu manifold}{Wu 流形}{〔保留〕}

\termgroup{六、K 理论与配边}
\termrowB{K-theory}{K 理论}{写法「$K$ 理论」}
\termrowA{KO-theory}{KO 理论}
\termrowA{K-theoretic}{K 理论的}
\termrowA{equivariant K-theory}{等变 K 理论}
\termrowB{Chern character}{Chern 特征}{〔保留〕}
\termrowB{bordism}{配边}{「协边」弃用}
\termrowB{cobordism}{配边}{与 bordism 同译「配边」，不区分}
\termrowB{bordism group}{配边群}{记号 $\Omega_n(X)$、$\Omega_n^{\mathrm{spin}}$}
\termrowA{spin bordism}{自旋配边}
\termrowA{bordism class}{配边类}
\termrowA{boundary (of a manifold) 配边中的}{边界}
\termrowB{Thom spectrum}{Thom 谱}{〔保留〕}
\termrowA{spectrum}{谱}
\termrowB{$\Omega$-spectrum}{$\Omega$ 谱}{〔保留〕}
\termrowA{generalized homology}{广义同调}
\termrowA{generalized cohomology}{广义上同调}
\termrowA{reduced homology}{约化同调}
\termrowB{Brown representability}{Brown 可表示性}{〔保留〕}
\termrowB{Anderson dual}{Anderson 对偶}{〔保留〕 亦作 Pontryagin dual，正文按原文}
\termrowB{Pontryagin dual}{Pontryagin 对偶}{〔保留〕}
\termrowB{Atiyah--Hirzebruch--Leray--Serre (AHLS) spectral sequence}{AHLS 谱序列}{〔保留〕 首见写全称}
\termrowA{spectral sequence}{谱序列}
\termrowA{$E_2$ page}{$E_2$ 页}
\termrowA{filtration}{滤过}
\termrowA{differential (of a spectral sequence)}{微分}
\termrowA{convergence}{收敛}
\termrowA{extension problem}{扩张问题}
\termrowA{invertible phase}{可逆相}
\termrowB{SPT phase}{SPT 相}{〔保留〕 首见写「对称保护拓扑（SPT）相」}
\termrowA{symmetry protected topological phase}{对称保护拓扑相}
\termrowA{topological phase}{拓扑相}
\termrowA{topological order}{拓扑序}
\termrowA{topological quantum field theory}{拓扑量子场论}
\termrowB{invertible TQFT}{可逆 TQFT}{〔保留〕}

\termgroup{七、费米子、对称类与凝聚态}
\termrowA{fermion}{费米子}
\termrowA{boson}{玻色子}
\termrowB{Majorana fermion}{Majorana 费米子}{〔保留〕}
\termrowB{Dirac fermion}{Dirac 费米子}{〔保留〕}
\termrowB{Weyl fermion}{Weyl 费米子}{〔保留〕}
\termrowA{time reversal}{时间反演}
\termrowA{time-reversal invariant}{时间反演不变的}
\termrowA{particle--hole}{粒子--空穴}
\termrowA{charge conjugation}{电荷共轭}
\termrowA{symmetry class}{对称类}
\termrowA{tenfold way}{十重分类}
\termrowB{Bott periodicity}{Bott 周期性}{〔保留〕}
\termrowA{2-periodicity}{2-周期性}
\termrowA{8-periodicity}{8-周期性}
\termrowA{gapped}{有能隙的}
\termrowA{gapless}{无能隙的}
\termrowA{gapped Hamiltonian}{有能隙哈密顿量}
\termrowA{Hamiltonian}{哈密顿量}
\termrowA{ground state}{基态}
\termrowA{excitation}{激发}
\termrowB{bulk}{体}{bulk-boundary correspondence「体--边对应」}
\termrowA{edge mode}{边缘模}
\termrowA{surface state}{表面态}
\termrowA{band structure}{能带结构}
\termrowA{Brillouin zone}{布里渊区}
\termrowA{topological insulator}{拓扑绝缘体}
\termrowA{superconductor}{超导体}
\termrowA{superfluid}{超流体}
\termrowA{order parameter}{序参量}
\termrowB{lattice}{格点}{「定义在格点上的哈密顿量」；晶体语境作「晶格」}
\termrowA{crystal}{晶体}
\termrowA{spin chain}{自旋链}
\termrowA{qubit}{量子比特}
\termrowB{Ising model}{Ising 模型}{〔保留〕}
\termrowA{partition function}{配分函数}
\termrowA{path integral}{路径积分}
\termrowA{effective action}{有效作用量}
\termrowA{low-energy effective action}{低能有效作用量}
\termrowA{topological term}{拓扑项}
\termrowA{quantized}{量子化的}
\termrowA{quantization}{量子化}
\termrowA{compactification}{紧化}
\termrowA{moduli space}{模空间}
\termrowA{supersymmetry}{超对称}
\termrowA{string theory}{弦论}
\termrowA{conformal field theory}{共形场论}
\termrowA{modular function}{模函数}
\termrowA{density matrix}{密度矩阵}
\termrowA{Hilbert space}{希尔伯特空间}
\termrowB{Berry phase}{Berry 相}{〔保留〕}
\termrowA{wavefunction}{波函数}
\termrowA{group extension}{群扩张}
\termrowA{central extension}{中心扩张}
\termrowA{group cohomology}{群上同调}
\termrowA{free resolution}{自由消解}
\termrowA{projective resolution}{投射消解}
\termrowA{abstract group}{抽象群}
\termrowA{finite group}{有限群}
\termrowA{cyclic group}{循环群}
\termrowA{anomaly of $1+1$d systems}{$1+1$ 维系统的反常}
\termrowB{$d+1$ dimensions}{$d+1$ 维}{统一写「$d+1$ 维」}
\termrowA{number theory}{数论}
\termrowA{algebraic geometry}{代数几何}
\termrowA{mathematical logic}{数理逻辑}
\termrowA{set theory}{集合论}
\termrowB{Gödel's incompleteness theorem}{Gödel 不完备定理}{〔保留〕}
\termrowA{physics}{物理}

\termgroup{八、陈述环境与文本用语（统一译法）}
\termrowB{Definition}{定义}{灰框内「定义 2.1.」}
\termrowA{Theorem}{定理}
\termrowA{Lemma}{引理}
\termrowA{Proposition}{命题}
\termrowA{Corollary}{推论}
\termrowA{Example}{例}
\termrowA{Remark}{注}
\termrowA{Exercise}{练习}
\termrowB{Note}{注}{原书 Note 与 Remark 同译「注」时须保持编号可辨，正文照抄编号}
\termrowA{Question}{问}
\termrowA{Problem}{问题}
\termrowA{Claim}{断言}
\termrowA{Conjecture}{猜想}
\termrowA{Assumption}{假设}
\termrowA{Notation}{记号}
\termrowB{Proof}{证明}{行内「Proof.」作 \textbackslash{}textbf{证明.}}
\termrowB{Aside}{附注}{标题作「附注：……」}
\termrowA{Acknowledgements}{致谢}
\termrowA{References}{参考文献}
\termrowB{Section / Sec.}{§}{引用写「§2.1」，不用 \textbackslash{}ref}
\termrowB{Figure / Fig.}{图}{引用写「图 2.1」}
\termrowA{Table}{表}
\termrowB{equation (2.1)}{式 (2.1)}{引用写纯文本「式 (2.1)」}
\termrowA{in this section}{本节}
\termrowB{as we will see}{我们将会看到}{避免「正如我们将要看到的」冗长译法}
\termrowB{up to}{至多差一个 / 在……意义下}{按语境：up to isomorphism「在同构意义下」}
\termrowA{roughly speaking}{粗略地说}
\termrowA{for example}{例如}
\termrowA{in general}{一般说来}
\termrowA{on the other hand}{另一方面}
\termrowA{note that}{注意}
\termrowB{let us}{我们}{「让我们」弃用}
\termrowA{it is easy to see that}{容易看出}
\termrowA{one can show}{可以证明}
\termrowA{by definition}{按定义}
\termrowA{trivially}{平凡地}
\termrowA{vacuously}{空泛地}
\termrowA{apocryphal}{轶闻的}
\termrowB{co-pilot}{Co-pilot}{〔保留〕}

\termgroup{九、保留拉丁字母的人名与专名（禁止音译）}
\termrowB{Wigner}{Wigner}{〔保留〕}
\termrowB{Chern}{Chern}{〔保留〕}
\termrowB{Simons}{Simons}{〔保留〕}
\termrowB{Yang}{Yang}{〔保留〕}
\termrowB{Mills}{Mills}{〔保留〕}
\termrowB{Dirac}{Dirac}{〔保留〕}
\termrowB{Maxwell}{Maxwell}{〔保留〕}
\termrowB{Hodge}{Hodge}{〔保留〕}
\termrowB{de Rham}{de Rham}{〔保留〕}
\termrowB{Čech}{Čech}{〔保留〕}
\termrowB{Künneth}{Künneth}{〔保留〕}
\termrowB{Hurewicz}{Hurewicz}{〔保留〕}
\termrowB{Whitehead}{Whitehead}{〔保留〕}
\termrowB{Eilenberg}{Eilenberg}{〔保留〕}
\termrowB{MacLane}{MacLane}{〔保留〕}
\termrowB{Postnikov}{Postnikov}{〔保留〕}
\termrowB{Pontryagin}{Pontryagin}{〔保留〕}
\termrowB{Stiefel}{Stiefel}{〔保留〕}
\termrowB{Whitney}{Whitney}{〔保留〕}
\termrowB{Euler}{Euler}{〔保留〕}
\termrowB{Betti}{Betti}{〔保留〕}
\termrowB{Grassmann}{Grassmann}{〔保留〕}
\termrowB{Hopf}{Hopf}{〔保留〕}
\termrowB{Klein}{Klein}{〔保留〕}
\termrowB{Möbius}{Möbius}{〔保留〕}
\termrowB{Koszul}{Koszul}{〔保留〕}
\termrowB{Morita}{Morita}{〔保留〕}
\termrowB{Novikov}{Novikov}{〔保留〕}
\termrowB{Calabi}{Calabi}{〔保留〕}
\termrowB{Yau}{Yau}{〔保留〕}
\termrowB{Bott}{Bott}{〔保留〕}
\termrowB{Atiyah}{Atiyah}{〔保留〕}
\termrowB{Singer}{Singer}{〔保留〕}
\termrowB{Patodi}{Patodi}{〔保留〕}
\termrowB{Hirzebruch}{Hirzebruch}{〔保留〕}
\termrowB{Leray}{Leray}{〔保留〕}
\termrowB{Serre}{Serre}{〔保留〕}
\termrowB{Mayer}{Mayer}{〔保留〕}
\termrowB{Vietoris}{Vietoris}{〔保留〕}
\termrowB{Dixmier}{Dixmier}{〔保留〕}
\termrowB{Douady}{Douady}{〔保留〕}
\termrowB{Anderson}{Anderson}{〔保留〕}
\termrowB{Brown}{Brown}{〔保留〕}
\termrowB{Clifford}{Clifford}{〔保留〕}
\termrowB{Majorana}{Majorana}{〔保留〕}
\termrowB{Weyl}{Weyl}{〔保留〕}
\termrowB{Berry}{Berry}{〔保留〕}
\termrowB{Pfaffian}{Pfaffian}{〔保留〕}
\termrowB{Wu}{Wu}{〔保留〕}
\termrowB{Wess}{Wess}{〔保留〕}
\termrowB{Zumino}{Zumino}{〔保留〕}
\termrowB{Witten}{Witten}{〔保留〕}
\termrowB{Kitaev}{Kitaev}{〔保留〕}
\termrowB{Thouless}{Thouless}{〔保留〕}
\termrowB{Kohmoto}{Kohmoto}{〔保留〕}
\termrowB{Nightingale}{Nightingale}{〔保留〕}
\termrowB{den Nijs}{den Nijs}{〔保留〕}
\termrowB{Haldane}{Haldane}{〔保留〕}
\termrowB{Kane}{Kane}{〔保留〕}
\termrowB{Mele}{Mele}{〔保留〕}
\termrowB{Fu}{Fu}{〔保留〕}
\termrowB{Kane--Mele}{Kane--Mele}{〔保留〕}
\termrowB{Schnyder}{Schnyder}{〔保留〕}
\termrowB{Ryu}{Ryu}{〔保留〕}
\termrowB{Furusaki}{Furusaki}{〔保留〕}
\termrowB{Ludwig}{Ludwig}{〔保留〕}
\termrowB{Altland}{Altland}{〔保留〕}
\termrowB{Zirnbauer}{Zirnbauer}{〔保留〕}
\termrowB{Teo}{Teo}{〔保留〕}
\termrowB{Kane}{Kane}{〔保留〕}
\termrowB{Mieda (三枝)}{Mieda}{〔保留〕 原文引用键 [三24]，正文保留}
\termrowB{Thouless}{Thouless}{〔保留〕}

\termgroup{十、易错与禁用清单}
\termrowB{division algebra}{可除代数}{标准译名；来源 ch02}
\termrowB{desingularization / desingularize}{消奇点}{代数几何通行译名；原句 We can desingularize $X=\C^2/\Z_2$；来源 ch02}
\termrowB{link（of a vertex）}{链环}{PL 拓扑标准译名；“the link of v”＝「v 的链环」；来源 ch02}
\termrowB{exotic sphere / exotic smooth structure}{异种球面 / 异种光滑结构}{微分拓扑通行译名；来源 ch02}
\termrowB{unbroken subgroup}{未破缺子群}{与「对称性破缺」配套；来源 ch02}
\termrowB{diagonal subgroup}{对角子群}{$G_{\text{diagonal}}=\{(g,g)\}$；来源 ch02}
\termrowB{phase diagram}{相图}{凝聚态通行译名；来源 ch02}
\termrowB{Cooper pair}{库珀对}{人名按第 1 条保留 Cooper？已处理：中文「库珀对」为通行定名，保留中文；来源 ch02}
\termrowB{degeneracy / degenerate}{简并 / 简并的}{量子力学通行译名；来源 ch02}
\termrowB{spin-orbit coupling}{自旋-轨道耦合}{原文 “spin-orbit coupling”；来源 ch02}
\termrowB{diffeomorphism}{微分同胚}{术语表仅收 homeomorphism 同胚，未收此词；来源 ch02}
\termrowB{piecewise linear (PL)}{分段线性（PL）}{原文以 PL 缩写为主，译文首次出现写全称后直接用 PL；来源 ch02}
\termrowB{norm-preserving}{保范的}{原文 “operators preserving the norm” 亦译「保范算子」；来源 ch02}
\termrowB{expectation value}{期望值}{序参量为 Cooper 对期望值；来源 ch02}
\termrowB{orbifold}{轨形}{存疑已裁决：轨形（不用「轨形空间」）；首见夹注 orbifold；来源 ch02}
\termrowB{multi-donut}{多甜甜圈}{存疑已裁决：多甜甜圈；首见给英文 multi-donut；来源 ch02}
\termrowB{crosscap}{交叉帽}{存疑已裁决：交叉帽（原文夹注「叉帽」，首见可注明）；后续统一「交叉帽」；来源 ch02}
\termrowB{texture}{织构}{原文括注 "texture"，指无核的拓扑孤子；先写「织构（texture）」再单用「织构」；来源 ch0506}
\termrowB{fine structure constant}{精细结构常数}{$\alpha=e^{2}/(4\pi\varepsilon_{0}\hbar c)\approx 1/137$；来源 ch0506}
\termrowB{Chern number}{Chern 数}{术语表只有 Chern class（Chern 类），未收 number；来源 ch0506}
\termrowB{magnetic flux}{磁通量}{术语表未收；与 Chern 数配对出现；来源 ch0506}
\termrowB{vector potential}{向量势}{术语表只有 gauge potential（规范势）；A 一律译「向量势」；来源 ch0506}
\termrowB{top form}{顶形式}{$\wedge^nT^*M$ 的截面；来源 ch0506}
\termrowB{density}{密度}{(5.27)–(5.29) 的 $\rho\,d^n x$ 型对象；与 probability density 区分；来源 ch0506}
\termrowB{inclusion map}{包含映射}{$\iota:N\hookrightarrow M$；来源 ch0506}
\termrowB{Hodge star}{Hodge 星号}{$\ast$ 算子；人名保留；来源 ch0506}
\termrowB{Dirac quantization condition}{Dirac 量子化条件}{〔保留〕 人名保留；来源 ch0506}
\termrowB{’t Hooft–Polyakov monopole}{’t Hooft–Polyakov 磁单极子}{〔保留〕 人名保留；来源 ch0506}
\termrowB{Chern–Simons term}{Chern–Simons 项}{〔保留〕 人名保留；连字符用 en dash；来源 ch0506}
\termrowB{singular (co)homology}{奇异（上）同调}{术语表只有 homology/cohomology；来源 ch0506}
\termrowB{simplicial complex}{单纯复形}{abstract simplicial complex 抽象单纯复形；来源 ch0506}
\termrowB{cell / cellular (co)homology}{胞腔 / 胞腔（上）同调}{n-cell 译「$n$-胞腔」；来源 ch0506}
\termrowB{attaching map}{贴附映射}{$f_i:\partial e_i\to X'$；来源 ch0506}
\termrowB{lens space}{透镜空间}{$S^{2n-1}/\Z_k$；来源 ch0506}
\termrowB{graded commutative}{分次交换}{$\omega\wedge\eta=(-1)^{pq}\eta\wedge\omega$；来源 ch0506}
\termrowB{Hopf invariant}{Hopf 不变量}{〔保留〕 人名保留；来源 ch0506}
\termrowB{perfect pairing}{完美配对}{值域为整个 $\Z$ 的配对；来源 ch0506}
\termrowB{Hurewicz homomorphism / theorem}{Hurewicz 同态 / Hurewicz 定理}{〔保留〕 人名保留；来源 ch0506}
\termrowB{homology sphere}{同调球}{与 $S^3$ 同调群相同的三维流形；来源 ch0506}
\termrowB{binary icosahedral group}{二元二十面体群}{$\hat\Gamma\subset SU(2)$；来源 ch0506}
\termrowB{reflection reality}{反射实性}{原文 reflection reality condition，指由 $\sigma$ 作用后再取复共轭的实性条件；来源 ch12}
\termrowB{reflection positivity}{反射正性}{来源 ch12}
\termrowB{null bordism}{零配边}{a null bordism：以给定流形为边界的那个配边；来源 ch12}
\termrowB{bordant}{配边的}{two manifolds are bordant →「两个流形配边」；来源 ch12}
\termrowB{incoming / outgoing boundary}{入边界 / 出边界}{配边 N 的两个边界分量；来源 ch12}
\termrowB{reduced bordism group}{约化配边群}{$\wt\Omega^{\mathcal S}_d(X)$；术语表作「约化配边」；来源 ch12}
\termrowB{torsion part}{挠部分}{术语表有「挠子群」，本章正文用「挠部分」；来源 ch12}
\termrowB{non-torsion part}{非挠部分}{来源 ch12}
\termrowB{A-roof genus ($\wh A$-genus)}{Â 亏格}{排为 $\wh{A}$ 亏格；来源 ch12}
\termrowB{Hermitean matrix}{Hermite 矩阵}{原文拼作 Hermitean，统一为 Hermite；来源 ch12}
\termrowB{thermal quantum Hall effect}{热量子霍尔效应}{来源 ch12}
\termrowB{Kramers degeneracy}{Kramers 简并}{来源 ch12}
\termrowB{quaternionic structure}{四元数结构}{来源 ch12}
\termrowB{zero mode}{零模}{亦作零本征值；来源 ch12}
\termrowB{chirality operator}{手征算子}{$\Gamma$，原文 the chirality operator；来源 ch12}
\termrowB{Arf invariant}{Arf 不变量}{来源 ch12}
\termrowB{Kitaev chain}{Kitaev 链}{来源 ch12}
\termrowB{Atiyah--Hirzebruch spectral sequence}{Atiyah--Hirzebruch 谱序列}{来源 ch12}
\termrowB{Adams spectral sequence}{Adams 谱序列}{来源 ch12}
\termrowB{deformation class}{形变类}{原文 deformation classes；来源 ch12}
\termrowB{semi ad-hoc analysis}{半凑合的分析}{原文 a semi ad-hoc analysis，保留作者自嘲语气；来源 ch12}
\termrowB{large gauge transformation}{大规范变换}{来源 ch12}
\termrowB{antipodal identification}{对径认同}{$\RP^2,\RP^3$ 上的 $x\sim-x$；来源 ch12}
\termrowB{pinor / (s)pinor}{pinor / (s)pinor}{〔保留〕 pin 结构的旋量，与 spinor（旋量）区分；来源 ch12}
\termrowB{NS / R boundary condition}{NS（反周期）/ R（周期）边界条件}{$S^1$ 上费米子的两种边界条件；来源 ch12}
\termrowB{periodic torus}{周期环面}{the periodic torus $T^2$；来源 ch12}
\termrowB{closed manifold}{闭流形}{无边紧流形；来源 ch12}
\termrowB{gluing}{粘合}{来源 ch12}
\termrowB{orientation reversal}{定向反转}{$\ol{M}$ 记号；来源 ch12}
\termrowB{division algebra}{可除代数}{标准译名；来源 ch02}
\termrowB{orbifold}{轨形}{存疑已裁决：轨形（不用「轨形空间」）；首见夹注 orbifold；来源 ch02}
\termrowB{multi-donut}{多甜甜圈}{存疑已裁决：多甜甜圈；首见给英文 multi-donut；来源 ch02}
\termrowB{crosscap}{交叉帽}{存疑已裁决：交叉帽（原文夹注「叉帽」，首见可注明）；后续统一「交叉帽」；来源 ch02}
\termrowB{Fact}{事实}{原书 Fact 灰框；preamble 无 fact 环境，用 `qmbox` + `\textbackslash{}qmstmt{事实}{n}` 排版；来源 ch03}
\termrowB{Riemannian metric}{Riemann 度量}{人名保留 Riemann；与已有 Riemannian manifold 体例一致；来源 ch03}
\termrowB{Lorentzian metric}{Lorentz 度量}{人名保留 Lorentz；来源 ch03}
\termrowB{orthonormal basis / frame}{正交归一基}{标准译名；Euclid 情形用「正交归一」，Lorentz 情形另注 $\eta_{ab}$；来源 ch03}
\termrowB{sphere bundle}{球丛}{词表仅有 unit sphere bundle 之外的纤维丛词条，此处登记；来源 ch03}
\termrowB{unit sphere bundle}{单位球丛}{沿用术语表 `Sph(E)` 记号；来源 ch03}
\termrowB{complex conjugate bundle}{复共轭丛}{与已有 conjugate 体例一致；来源 ch03}
\termrowB{direct sum (of bundles)}{直和}{逐纤维直和；来源 ch03}
\termrowB{tensor product (of bundles)}{张量积}{逐纤维张量积；来源 ch03}
\termrowB{double cover}{双重覆叠}{原书 “double cover of $SO(3)$”；覆叠空间标准译名；来源 ch03}
\termrowB{spinor bundle}{旋量丛}{与已有 spinor 旋量体例一致；来源 ch03}
\termrowB{spinor representation}{旋量表示}{群表示论标准译名；来源 ch03}
\termrowB{octonion}{八元数}{标准译名；原书 $S^7$ 与八元数乘法；来源 ch03}
\termrowB{Bloch sphere}{Bloch 球面}{人名保留 Bloch；首见「Bloch 球面」；来源 ch03}
\termrowB{tetrad / vierbein / fielbein}{tetrad / vierbein / fielbein（保留原词）}{原书三个同义词并列表述，译文保留原词并夹注「标架」类说明；来源 ch03}
\termrowB{electroweak theory}{电弱理论}{粒子物理标准译名；来源 ch03}
\termrowB{symmetry breaking}{对称破缺}{与已有 symmetry 体例一致；原文 “pattern of symmetry breaking” 译「对称破缺模式」；来源 ch03}
\termrowB{Higgs boson}{Higgs 玻色子}{人名保留 Higgs，与 Higgs field 一致；来源 ch03}
\termrowB{quark}{夸克}{通行译名；来源 ch03}
\termrowB{lepton}{轻子}{通行译名；来源 ch03}
\termrowB{hypercharge}{超荷}{粒子物理标准译名（$U(1)$ 超荷）；来源 ch03}
\termrowB{generation}{代}{原书仅用于下标 $V_{\mathrm{single\ generation}}$，译文保留英文下标，未在正文另译；来源 ch03}
\termrowB{induced homomorphism}{诱导同态}{同伦群函子性标准译名；来源 ch04}
\termrowB{n-connected}{$n$-连通}{与 connected 连通、simply connected 单连通体例一致；来源 ch04}
\termrowB{generic point}{一般位置的点}{原文 “a point $p$ in general position”；亦称「一般点」，正文统一用「一般位置的点」；来源 ch04}
\termrowB{presentation}{展示}{群论标准译名（generators and relations）；来源 ch04}
\termrowB{commutator}{交换子}{群论标准译名；来源 ch04}
\termrowB{linking number}{环绕数}{与已有 winding number 缠绕数并列；$\mathrm{Lk}$ 记号保留；来源 ch04}
\termrowB{Hopf fibration}{Hopf 纤维化}{保留人名；$S^1\to S^3\to S^2$ 等三类；来源 ch04}
\termrowB{reduced suspension}{约化悬挂}{与 suspension 悬挂（已收）区分；记号 $\Sigma X$；来源 ch04}
\termrowB{suspension homomorphism}{悬挂同态}{与 suspension isomorphism 悬挂同构（已收）一致；来源 ch04}
\termrowB{stable homotopy group}{稳定同伦群}{灰色区域内的球面同伦群；来源 ch04}
\termrowB{unstable group}{不稳定群}{非灰色区域；来源 ch04}
\termrowB{stem}{stem（保留）}{原书 “the $k$-th stem”，学界通用英文术语，正文保留并夹注「第 $k$ 个 stem」；来源 ch04}
\termrowB{Freudenthal suspension theorem}{Freudenthal 悬挂定理}{人名保留 Freudenthal；来源 ch04}
\termrowB{Mahowald uncertainty principle}{Mahowald 不确定性原理}{人名保留 Mahowald；来源 ch04}
\termrowB{extension (of groups)}{扩张}{原书 “$B$ is an extension of $C$ by $A$」，译「$B$ 为 $C$ 被 $A$ 的扩张」；来源 ch04}
\termrowB{exactness}{正合性}{与 exact sequence 正合列一致；来源 ch04}
\termrowB{trivial bundle}{平凡丛}{与 local trivialization 局部平凡化一致；来源 ch04}
\termrowB{compact Lie group}{紧李群}{「紧」不写「紧致」（术语表错误项）；来源 ch04}
\termrowB{simple Lie group}{单李群}{与 semisimple 半单对应；来源 ch04}
\termrowB{spin bundle}{旋量丛}{与 ch03 spinor bundle 旋量丛统一；来源 ch04}
\termrowB{associated vector bundle}{相伴向量丛}{associated bundle 相伴丛（已收）+ vector bundle 向量丛（已收）；来源 ch04}
\termrowB{instanton number}{瞬子数}{与 instanton 瞬子（已收）一致；来源 ch04}
\termrowB{instanton particle}{瞬子粒子}{原书 “Instanton particles”；来源 ch04}
\termrowB{magnetic charge}{磁荷}{与 electric charge 电荷对应；来源 ch04}
\termrowB{magnetic monopole}{磁单极子}{术语表已收 monopole 磁单极子；来源 ch04}
\termrowB{mod-2 anomaly}{模 2 反常}{Witten 的 $SU(2)$ 模 2 反常；来源 ch04}
\termrowB{field configuration}{场组态}{原书 “field configuration”；来源 ch04}
\termrowB{vortex}{涡旋}{复数 vortices；Abrikosov 涡旋；来源 ch04}
\termrowB{core}{芯}{原书 “without a core / with a core」，译「无芯 / 有芯」；来源 ch04}
\termrowB{codimension}{余维数}{原书 “codimension $d$”；来源 ch04}
\termrowB{topological charge}{拓扑荷}{与 topological soliton 拓扑孤子（已收）一致；来源 ch04}
\termrowB{non-invertible symmetry}{非可逆对称性}{原书 “non-invertible symmetry”，近年凝聚态/高能常用；来源 ch04}
\termrowB{baby Skyrmion}{baby Skyrmion（保留）}{原书术语，保留英文并夹注「常简称为 Skyrmion」；来源 ch04}
\termrowB{Hopfion}{Hopfion（保留）}{由 Hopf 不变量得名，学界通用，保留英文；来源 ch04}
\termrowB{D-brane}{D-膜}{原书 “D0/D4-branes”；来源 ch04}
\termrowB{'t Hooft–Polyakov monopole}{'t Hooft–Polyakov 磁单极子}{人名保留；来源 ch04}
\termrowB{Abrikosov–Nielsen–Olesen vortex}{Abrikosov–Nielsen–Olesen 涡旋}{人名保留；原书印 “Nielsen-Olsen”，按通行拼法 Olesen；来源 ch04}
\termrowB{good cover（足够细的覆盖）}{足够细的覆盖}{原文用 “sufficiently fine cover”；未用 good cover 之名；来源 ch0708}
\termrowB{nerve}{神经}{原文括注 nerve；「覆盖的神经」；来源 ch0708}
\termrowB{Čech p-cochain}{Čech $p$-上链}{与 singular cochain 区分时加前缀；来源 ch0708}
\termrowB{holonomy}{和乐}{原文「对数学的人，叫和乐（holonomy）」；物理语境可理解为 Wilson 圈的相位；来源 ch0708}
\termrowB{pasting function}{粘贴函数}{$g_{ij}(x)$；与 transition function 同指；来源 ch0708}
\termrowB{trivialization}{平凡化}{「丛平凡化的障碍」；来源 ch0708}
\termrowB{obstruction}{障碍}{术语表未收；obstruction class 障碍类；来源 ch0708}
\termrowB{reduction of structure group}{结构群的约化}{原文 “reduce the structure group”；来源 ch0708}
\termrowB{Stiefel–Whitney class}{Stiefel--Whitney 类}{〔保留〕 $w_1,w_2,w_3$；「第 1/2/3 Stiefel--Whitney 类」；来源 ch0708}
\termrowB{integral Stiefel–Whitney class}{整 Stiefel--Whitney 类}{$W_3(P)$；原文 “integral third Stiefel–Whitney class”；来源 ch0708}
\termrowB{Bockstein homomorphism}{Bockstein 同态}{〔保留〕 $\beta:H^2(M;B)\to H^3(M;A)$；人名不音译；来源 ch0708}
\termrowB{Eilenberg–Mac Lane space}{Eilenberg--Mac Lane 空间}{〔保留〕 $K(A,k)$；人名保留；来源 ch0708}
\termrowB{mod 2 reduction}{模 2 约化}{$c_1(L)\bmod 2$；全章多处；来源 ch0708}
\termrowB{integer quantum Hall effect}{整数量子霍尔效应}{术语表未收；IQHE；来源 ch0708}
\termrowB{fractional quantum Hall effect}{分数量子霍尔效应}{术语表未收；FQHE；来源 ch0708}
\termrowB{conductance tensor}{电导张量}{(8.18) 的 $\sigma^{\text{SI}}$；对称部分与反对称部分；来源 ch0708}
\termrowB{gap / gapped}{能隙 / 有能隙}{表 4 表头「有能隙？」；「无能隙」不量子化；来源 ch0708}
\termrowB{abelian Chern–Simons term}{阿贝尔 Chern--Simons 项}{与「非阿贝尔」相对；来源 ch0708}
\termrowB{gravitational Chern–Simons term}{引力 Chern--Simons 项}{$\mathrm{CS}(\omega)$；来源 ch0708}
\termrowB{thermal Hall conductance}{热霍尔电导}{与电荷霍尔电导对照；来源 ch0708}
\termrowB{chiral fermion}{手征费米子}{“chiral fermions in four dimensions”；来源 ch0708}
\termrowB{Steenrod square}{Steenrod 平方}{〔保留〕 $Sq^k$；人名保留；来源 ch0708}
\termrowB{Wu class}{Wu 类}{〔保留〕 $\nu_k$；人名保留，不音译；来源 ch0708}
\termrowB{quadratic refinement}{二次精化}{原文括注 quadratic refinement；指把歧义减半的精细化；来源 ch0708}
\termrowB{dual triangulation}{对偶三角剖分}{畴壁即对偶图的边；来源 ch0910}
\termrowB{domain wall}{畴壁}{原文 “topological $G$ domain wall”；术语表未收，全章高频；来源 ch0910}
\termrowB{Dijkgraaf–Witten theory}{Dijkgraaf--Witten 理论}{〔保留〕 [DW90]；有限群规范理论；来源 ch0910}
\termrowB{Dijkgraaf–Witten phase}{Dijkgraaf--Witten 相}{〔保留〕 可逆拓扑相；来源 ch0910}
\termrowB{projective representation}{射影表示}{$U(g_1)U(g_2)=\alpha(g_1,g_2)U(g_1g_2)$；来源 ch0910}
\termrowB{projective phase}{射影相位}{$\alpha(g_1,g_2)\in U(1)$；来源 ch0910}
\termrowB{sporadic finite simple group}{散在有限单群}{术语表未收；来源 ch0910}
\termrowB{monster group}{monster 群}{〔保留〕 原文小写 m；最大散在有限单群，$M$；来源 ch0910}
\termrowB{GAP / SAGE / Magma}{GAP / SAGE / Magma}{〔保留〕 计算机代数系统，不译；来源 ch0910}
\termrowB{complex fermion}{复费米子}{$c_i$，$\{c_i,c_j^\dagger\}=\delta_{ij}$；来源 ch0910}
\termrowB{real fermion}{实费米子}{与 Majorana fermion 同义；来源 ch0910}
\termrowB{fermion parity}{费米子宇称}{$(-1)^F$；与「宇称反常」的宇称同字；来源 ch0910}
\termrowB{ungraded representation}{无分次表示}{Clifford 代数的通常表示；来源 ch0910}
\termrowB{graded representation}{分次表示}{带 $(-1)^F$ 的表示；术语表有「分次交换」；来源 ch0910}
\termrowB{real Clifford algebra}{实 Clifford 代数}{$\mathrm{Cliff}_{\R}(n_+,n_-)$；来源 ch0910}
\termrowB{Arf theory}{Arf 理论}{Kitaev 链非平凡相的别称；术语表有 Arf 不变量；来源 ch0910}
\termrowB{time-reversal invariance}{时间反演不变性}{$T^2=+1$，$T$ 反酉；来源 ch0910}
\termrowB{antiunitary}{反酉的}{$Ti=-iT$；来源 ch0910}
\termrowB{complex conjugation}{复共轭}{$T$ 取 $z_1,\dots,z_N$ 的复共轭；来源 ch0910}
\termrowB{Hermitian}{Hermite 的}{「Hermite 组合」「反 Hermite 矩阵」；术语表有「Hermite 矩阵」，全书不用「厄米」；来源 ch0910}
\termrowB{Morita equivalence}{Morita 等价}{术语表有 Morita（保留）；来源 ch0910}
\termrowB{representation category}{表示范畴}{$\Rep(A)$；来源 ch0910}
\termrowB{category equivalence}{范畴等价}{Morita 等价即表示范畴等价；来源 ch0910}
\termrowB{matrix algebra}{矩阵代数}{$M_2(\R)$、$M_2(\HP)$；来源 ch0910}
\termrowB{Dirac equation}{Dirac 方程}{$(\gamma^\mu\partial_\mu+m)\psi=0$；来源 ch0910}
\termrowB{Lorentz generator}{Lorentz 生成元}{〔保留〕 $\gamma^\mu\gamma^\nu$；来源 ch0910}
\termrowB{flavor symmetry}{味对称}{复数情形 $u(n)$，实数情形 $so(n)/u(n)/sp(n)$；来源 ch0910}
\termrowB{reflection}{反射}{关于超平面的反射；来源 ch0910}
\termrowB{hyperplane}{超平面}{与 $w=(w^a)$ 垂直的超平面；来源 ch0910}
\termrowB{kernel}{核}{$\Spin(n)\to\SO(n)$ 的核为 $\{\pm1\}$；来源 ch0910}
\termrowB{gapping term}{开能隙项}{原文 “a mass term or a gapping term”；来源 ch0910}
\termrowB{gravitational Chern--Simons term}{引力 Chern--Simons 项}{2+1 维体低能有效作用量中的该拓扑项，给出整数热霍尔电导；来源 ch1113}
\termrowB{gapped mass term}{有能隙质量项}{表 7/表 8 列名 gapped Hamiltonians 的等价说法；来源 ch1113}
\termrowB{essentially gapless degrees of freedom}{本质上无能隙的自由度}{原书加引号强调：可加质量的自由度不算；来源 ch1113}
\termrowB{symmetric space}{对称空间}{齐性空间 $G/H$ 上带 2 阶自同构 $\sigma$ 者；来源 ch1113}
\termrowB{Cartan classification}{Cartan 分类}{十族对称空间；正文并出现「Cartan 名（A、AIII、AI…）」；来源 ch1113}
\termrowB{random matrix theory}{随机矩阵理论}{Wigner--Dyson 三类系综（GUE/GOE/GSE）的来源；来源 ch1113}
\termrowB{Bogoliubov--de Gennes (BdG) class}{BdG 类}{原书 'BdG' or 'AZ' symmetry classes；AZ 指 Altland--Zirnbauer；来源 ch1113}
\termrowB{particle--hole symmetry}{粒子--空穴对称}{反幺正 $\tilde T$，$\tilde TH=-H\tilde T$；来源 ch1113}
\termrowB{chiral symmetry}{手征对称性}{存在幺正 $S$ 使 $SH+HS=0$；与「手征性」区分；来源 ch1113}
\termrowB{classifying space of invertible phases}{可逆相的分类空间}{$\mathcal H_d$；Kitaev 论证的出发点；来源 ch1113}
\termrowB{dimension axiom}{维数公理}{Eilenberg--Steenrod 第 (5) 条公理；来源 ch1113}
\termrowB{Eilenberg--Steenrod axioms}{Eilenberg--Steenrod 公理}{普通（上）同调的公理化刻画；来源 ch1113}
\termrowB{generalized cohomology theory}{广义上同调理论}{术语表收「广义（上）同调」，此处按原文 co 前缀分写；来源 ch1113}
\termrowB{associated graded}{相伴分次}{$\operatorname{Gr}(T_d(X))$，谱序列极限给出的对象；来源 ch1113}
\termrowB{homotopy long exact sequence}{同伦长正合列}{纤维化的 $\pi_*$ 长正合列；来源 ch1113}
\termrowB{Leray--Serre spectral sequence}{Leray--Serre 谱序列}{AHLS 谱序列在 $T=H$ 时的特例；来源 ch1113}
\termrowB{bordism cohomology}{配边上同调}{$\Omega^d_{\mathcal S}(X)$；原书又称 cobordism group；来源 ch1113}
\termrowB{proper map}{proper 映射}{式 (13.21) 中余维 $d$ 的 proper 映射（逆紧映射）；来源 ch1113}
\termrowB{cohomology operation}{上同调运算}{Bockstein 与 $\mathcal P^1$ 均属此类；来源 ch1113}
\termrowB{Bockstein (homomorphism)}{Bockstein（同态）}{$\bar\beta:H^d(X;\Z_3)\to H^{d+1}(X;\Z)$；系数列 $0\to\R/\Z\xrightarrow{3\times}\R/\Z\to\Z_3\to0$；来源 ch1113}
\termrowB{skeleton (of a CW complex)}{骨架}{$B_0\subset B_1\subset\cdots$（术语表已收「骨架」，此处标明确切语境）；来源 ch1113}
\termrowB{spectral sequence page}{谱序列的页}{$E^2$ 页、$E^3$ 页、$E_r$ 页、$E_\infty$ 页；上标/下标两种记法照原书；来源 ch1113}
\termrowB{differential of a spectral sequence}{谱序列的微分}{$d^r$（同调）与 $d_r$（上同调）两种排法，照原书不统一；来源 ch1113}
\termrowB{unitary}{酉}{术语表原缺（Lead m00827 定：取「酉」，不用「幺正」）；ch11 原 44 处「幺正」已统一改，全书现 66 处「酉」；来源 consistency}
\termrowB{operator}{算子}{术语表原缺；原「算符」4 处已改为「算子」，全书现 108 处；来源 consistency}
\termrowB{universal covering}{万有覆叠}{术语表原缺；另有「万有覆叠空间」写法，指同一对象，保留；来源 consistency}
\end{multicols}
